\chapter{Project Timeline (Gantt Chart)}
\label{app:gantt}

This appendix presents the project timeline as a Gantt-style overview.

%==============================================================================
\section{12-Month Project Timeline}
%==============================================================================

The project follows a 12-month schedule aligned with the proposal. Phases 1--9 span literature review, cohort design, tokenization, ETHOS implementation, few-shot and federated training, baseline comparison, experiments, and thesis writing.

%==============================================================================
\section{Phase Overview}
%==============================================================================

\begin{table}[h]
\centering
\caption{Project Gantt Overview}
\begin{tabular}{llrr}
\toprule
Phase & Task & Start (Week) & Duration \\
\midrule
1 & Literature review & 1 & 4 \\
2 & MIMIC-IV access \& cohort design & 5 & 2 \\
3 & Data extraction \& tokenization & 7 & 4 \\
4 & ETHOS encoder implementation & 11 & 3 \\
5 & Few-shot meta-learning & 14 & 4 \\
6 & Federated training & 18 & 2 \\
7 & Tabular SOTA baseline & 12 & 4 \\
8 & Experiments \& ablation & 20 & 4 \\
9 & Writing \& figures & 24 & 8 \\
\bottomrule
\end{tabular}
\label{tab:gantt}
\end{table}

%==============================================================================
\section{Key Milestones}
%==============================================================================

\begin{itemize}
\item \textbf{Week 4}: Literature review complete; research questions finalized.
\item \textbf{Week 7}: Cohort extracted; first 2,000 patients processed.
\item \textbf{Week 14}: ETHOS+FewShot pipeline operational.
\item \textbf{Week 18}: Federated vs.\ centralized comparison complete.
\item \textbf{Week 24}: All experiments run; best model (RF) identified.
\item \textbf{Week 32}: Thesis draft complete; defense preparation.
\end{itemize}

%==============================================================================
\section{Roadmap Graphic}
%==============================================================================

Figure~\ref{fig:roadmap} illustrates the future work roadmap (1--5 years).

\begin{figure}[htbp]
\centering
\fbox{\parbox{0.8\textwidth}{
\centering
\textbf{Future Work Roadmap}\\[0.5em]
Phase 1 (1 yr): Larger cohort, external validation\\
Phase 2 (2 yr): Multi-site federated trial\\
Phase 3 (3 yr): Prospective deployment\\
Phase 4 (4 yr): Regulatory approval path\\
Phase 5 (5 yr): Clinical decision support integration
}}
\caption{Future work phases (1--5 years)}
\label{fig:roadmap}
\end{figure}

%==============================================================================
\section{Risk Management}
%==============================================================================

Table~\ref{tab:risk} summarizes project risks and mitigation strategies, following the proposal's project management framework.

\begin{table}[htbp]
\centering
\caption{Project Risk Management}
\begin{tabular}{llll}
\toprule
Risk & Probability & Impact & Mitigation \\
\midrule
Data access delay & Medium & High & Early MIMIC-IV credential application; backup synthetic data \\
Model underperformance & Low & Low & Achieved near-parity via stacking (0.732), GFA (0.737); classical baseline as fallback \\
Federated convergence issues & Low & Medium & FedAvg tuning; fallback to centralized \\
Computational resource limits & Low & Low & Cloud credits; reduced cohort if needed \\
Timeline slippage & Medium & Medium & Buffer in Gantt; prioritize core experiments \\
\bottomrule
\end{tabular}
\label{tab:risk}
\end{table}
